\inicializarSubseccion{Taller de interpolación – Alexa}

\begin{figure}[H]
    \centering
    \begin{minipage}[b]{0.30\textwidth}
        \centering
        \includegraphics[width=\textwidth]{\nat{alexaPuntosUno.png}}
        \caption{Puntos de la interpolación del taller (P1).}
        \label{fig:alexaPuntosUno}
    \end{minipage}
    \hfill
    \begin{minipage}[b]{0.30\textwidth}
        \centering
        \includegraphics[width=\textwidth]{\nat{alexaPuntosDos.png}}
        \caption{Puntos de la interpolación del taller (P2) y una función.}
        \label{fig:alexaPuntosDos}
    \end{minipage}
\end{figure}

\begin{figure}[H]
    \centering
    \includegraphics[width=0.30\textwidth]{\nat{alexaFunciones.png}}
    \caption{Los polinomios generados por la interpolación.}
    \label{fig:alexaFunciones}
\end{figure}

\begin{figure}[H]
    \centering
    \includegraphics[width=0.45\textwidth]{\nat{alexaPlanoGeogebra.png}}
    \caption{Gráfica generada por los puntos en Geogebra.}
    \label{fig:alexaPlanoGeogebra}
\end{figure}

\begin{figure}[H]
    \centering
    \includegraphics[width=0.45\textwidth]{\nat{alexaPlanoMATLAB.png}}
    \caption{Gráfica generada por los puntos en MATLAB.}
    \label{fig:alexaPlanoMATLAB}
\end{figure}